\documentclass[12pt]{article}
\usepackage{latexsym, amscd, %graphicx, 
amssymb}
\newcommand{\SECT}[1]
		{\vspace{.3cm}
		$\!\!\!\!\!\!\!\!\!\!\!${\bf {#1}}
		\vspace{.2cm}}

\begin{document}


\begin{center}

{\large \bf Math 152, Fall, 2004, Workshop 10\par}
\vspace{.5cm}
{{\bf Honors Section} 
\par}
\end{center}

\begin{enumerate}

\item \begin{enumerate}

\item For each power series below, use the ratio test to 
determine all values of $x$ for which the 
series converges absolutely, then analyze the behavior 
of the series at the endpoints in order to 
determine the interval of convergence. 
$$(i)\; \sum_{n=0}^{\infty}\frac{nx^{n}}{n^{2}+1},\;\;\;\;\;\;\;\;
(ii)\; \sum_{n=1}^{\infty}\frac{n^{2}(x-1)^{n}}{2^{n}},\;\;\;\;\;\;\;\;
(iii)\; \sum_{n=1}^{\infty}\frac{3^{n}x^{n}}{n^{2}}.$$

\item Can you cook up a power series whose interval of convergence 
is the interval $(0, 1]$, that is, the 
interval defined by $0 < x \le 1$? How about $(0, 1)$? 
Give an explicit series or explain why you can't.

\end{enumerate}

\item Consider the function defined by 
$$f(x) = \sum_{n=0}^{\infty}(-1)^{n}\frac{x^{2n}}{n!}.$$

\begin{enumerate}

\item Write out the first five terms of the series (up to $n = 4$; 
remember that $0! = 1$). 

\item Determine the interval of convergence; this is the domain of $f$. 

\item Verify in the following two ways that 
\begin{equation}
f'(x) = -2xf(x):
\end{equation}

\begin{enumerate}

\item Compute both sides of (1) in terms of the five terms 
you wrote out in; 

\item Verify (1) using the original summation (sigma) notation. This 
is a bit tricky (which is why it is 
a good idea to do i first); you will have to make a change 
of summation index. 

\end{enumerate}

\item Explain why $y = f(x)$ is a solution of the initial value problem 
$$y' = 2xy,\;\;\;\;\;\;y(0) = 1.$$

\item Solve this initial value problem and thereby obtain a formula 
for f(x) in terms of 
"elementary functions" (those found on your calculator). 

\end{enumerate}

\item Keep the notation of the previous problem, and let 
$$s_{N}(x) =\sum_{n=1}^{N}(-1)^{n}\frac{x^{2n}}{n!},$$ 
e.g., $s_{3}(x) = 1- x^{2} + x^{4}/2- x^{6}/6$. 

\begin{enumerate}

\item Using the formula you discovered for $f(x)$, use your calculator 
to graph $f$ and the 
partial sums $s_{0}$, $s_{2}$, $s_{4}$, and $s_{6}$ at the same time, 
in a window where $0 \le x \le 1.2$. Copy 
your graph to your workshop, identifying the different curves. 
Does it appear that $s_{N}$ 
becomes a better approximation to $f$ as $N$ increases? 

\item Use the alternating series error formula to obtain an 
upper bound for the error in the 
approximation $f(x) \simeq s_{6}(x)$, $0 \le x \le 1.2$. 
Your answer should be a single number that 
applies to all $x$ in the range $0 \le x \le 1.2$. Explain why 
this number is consistent with the 
graph from (a). 

\end{enumerate}

\end{enumerate}

\end{document}